% This file is part of the modular source for Stratified Relational Workspace Theory.
% !TeX root = ../Stratified_Relational_Workspace_Theory.tex
\theorypart{XVIII}{Specification Robustness, Identifiability, and Statistical Certification}

A workspace certificate is never evaluated in a vacuum. It is evaluated for a
particular future variable, loss, time horizon, intervention family, null
panel, threshold vector, and system decomposition. These choices determine
which distinctions matter and which capacities count as successful. They are
therefore part of the scientific claim rather than clerical details surrounding
it.

This dependence has two consequences. First, a result that passes under one
reasonable analysis can fail under another. Robustness must be assessed across
a declared family of scientifically defensible alternatives. Second, even a
fully specified claim may not be identifiable from the observations and
interventions that are actually available. Distinct realizations can induce the
same experimental law while disagreeing about whether a relational workspace
exists. The theory must say when such disagreement can be resolved, when only
bounds are justified, and how sampling uncertainty compounds these structural
ambiguities.

\topic{61}{The Complete Scientific Specification}

\begin{definition}[Workspace specification]\label{def:workspace-specification}
A workspace specification is a tuple
\[
\boxed{
\varsigma
=
\bigl(
\mathcal D,
Y^+,
\ell,
\tau,
\mathcal I,
\mathcal N,
\Theta
\bigr),
}
\]
where:
\begin{itemize}
\item \(\mathcal D\) is the decomposition and local-access model;
\item \(Y^+\) is the declared future variable or trajectory;
\item \(\ell\) is the normalized loss or utility convention;
\item \(\tau\) is the temporal horizon or horizon family;
\item \(\mathcal I\) is the feasible intervention family;
\item \(\mathcal N\) is the required matched-null family; and
\item \(\Theta\) is the gate schema, direction, threshold, and indifference
 region used for certification.
\end{itemize}
The complete certificate is written
\[
\operatorname{Cert}(W;\varsigma)
=
\bigwedge_{q\in\mathcal J(\varsigma)}
\mathbf 1\{m_q(W;\varsigma)>0\},
\]
where \(m_q\) is the signed margin of gate \(q\). Positive margin means that
that gate passes in the direction declared by \(\Theta\).
\end{definition}

The notation absorbs lower-bounded and upper-bounded gates into one sign
convention. If a causal effect must exceed \(c_0\), then
\[
m_q=\widehat\eta-c_0.
\]
If a defect must remain below \(\varepsilon_0\), then
\[
m_q=\varepsilon_0-\widehat\delta.
\]
The complete certificate therefore passes exactly when every required signed
margin is positive.

The specification does not include arbitrary features of the analysis. It
contains choices that alter the meaning or evidential force of the claim. An
optimization algorithm, numerical tolerance, or estimator belongs in the
estimation protocol unless changing it changes the target estimand or the
admissible interventions. Conversely, the target, horizon, and null family
cannot be treated as implementation details because changing them can reverse
the conclusion without changing the underlying realization.

\subtopic{61.1 Specification relativity is not unrestricted relativism}

A theorem-development workspace may be operationally real for a month-long
research trajectory and irrelevant for a ten-second motor-control target. A
repair channel may preserve theorem identity while failing to preserve social
coordination. These are different scientific claims. The theory does not
require one workspace state to be sufficient for every possible future or loss.
It requires the intended scope to be declared before certification.

This yields the more precise notation
\[
\boxed{
\SRW(W;\varsigma)=1.
}
\]
The earlier expression \(\SRW_\Theta(W;\mathcal D)\) is the special case in
which the remaining coordinates of \(\varsigma\) are held fixed.

\subtopic{61.2 Inert reparameterization and substantive change}

Some changes of specification merely recode the same claim. Others change its
content. The distinction is analogous to gauge freedom in the structural
theory.

\begin{definition}[Specification equivalence]\label{def:spec-equivalence}
Two specifications \(\varsigma\) and \(\varsigma'\) are equivalent, written
\(\varsigma\simeq_{\mathrm{spec}}\varsigma'\), when there are bijective
transport maps that:
\begin{enumerate}
\item preserve the sigma-field generated by the future target;
\item preserve the ordering and normalization of decisions under the loss;
\item map feasible interventions to physically equivalent interventions;
\item map every required null to a null preserving the same quantities; and
\item carry every gate threshold and margin into its counterpart.
\end{enumerate}
\end{definition}

\begin{proposition}[Certificate invariance under specification equivalence]
\label{prop:spec-equivalence-invariance}
If \(\varsigma\simeq_{\mathrm{spec}}\varsigma'\) and \(\Phi\) is the induced
workspace correspondence, then
\[
\operatorname{Cert}(W;\varsigma)
=
\operatorname{Cert}(\Phi W;\varsigma').
\]
\end{proposition}

\begin{proof}
Every required gate is carried to a gate with the same signed margin. The
conjunction of positive margins is therefore unchanged.
\end{proof}

A bijective relabeling of target classes with the correspondingly transported
loss is inert. Replacing exact theorem validity with stylistic similarity is
not. Changing a one-step horizon to a year-long persistence horizon is not.
Removing the null that tests hidden centralization is not. Robustness concerns
substantive alternatives after equivalent recodings have been quotiented out.

\subtopic{61.3 Admissible specification families}

\begin{definition}[Admissible specification family]
\label{def:admissible-spec-family}
An admissible specification family \(\mathfrak S\) is a predeclared set of
specifications that:
\begin{enumerate}
\item refer to the same scientific claim type and carrier presentation;
\item use targets that are equivalent operationalizations or explicitly
 declared nearby stress targets;
\item normalize losses so that margins are comparable;
\item use horizons justified by the process being studied;
\item include every feasible intervention and null needed to distinguish the
 principal alternative explanations;
\item use thresholds fixed without access to certification outcomes; and
\item exclude only alternatives for reasons stated before evaluation.
\end{enumerate}
\end{definition}

The family should not contain every imaginable analysis. Such a requirement
would erase meaningful domain specificity. It should contain the alternatives
that a technically informed critic could reasonably regard as equally valid or
as necessary stress tests for the declared claim. Multiverse and specification-
curve methods formalize related concerns in empirical analysis
\cite{SteegenEtAl2016,SimonsohnEtAl2020}, while sensitivity analysis and
partial identification emphasize that conclusions should expose their
assumption dependence rather than conceal it \cite{Leamer1983,Manski2003}.

Three classes should be reported separately:
\begin{description}
\item[Equivalent alternatives.] Different encodings or estimators of the same
 estimand after appropriate transport.
\item[Robustness alternatives.] Scientifically defensible choices intended to
 support the same claim.
\item[Scope alternatives.] Genuinely different targets or domains that define
 different claims and should not be averaged into one robustness number.
\end{description}

\subtopic{61.4 Common gate schema}

A robustness family must use a common semantic gate index
\(\mathcal J_*\). A specification may mark a gate as not required only when the
claim level itself does not include that property. It may not remove a failed
gate after observing the result. Let
\[
\mathcal J(\varsigma)\subseteq\mathcal J_*
\]
be fixed before certification. For a family intended to support one fixed
claim level, the cleanest design uses
\[
\mathcal J(\varsigma)=\mathcal J_*
\qquad
(\varsigma\in\mathfrak S).
\]
This prevents apparent robustness from being created by silently changing the
meaning of passage across the family.

\topic{62}{Robustness Across Admissible Specifications}

For each gate \(q\in\mathcal J_*\), define the specification envelope
\[
\boxed{
\underline m_q^{\mathfrak S}(W)
=
\inf_{\varsigma\in\mathfrak S}m_q(W;\varsigma),
\qquad
\overline m_q^{\mathfrak S}(W)
=
\sup_{\varsigma\in\mathfrak S}m_q(W;\varsigma).
}
\]
Its sensitivity width is
\[
\Delta_q^{\mathrm{spec}}(W)
=
\overline m_q^{\mathfrak S}(W)
-
\underline m_q^{\mathfrak S}(W).
\]
The robust certificate margin is
\[
\boxed{
r_{\mathfrak S}(W)
=
\inf_{\varsigma\in\mathfrak S}
\min_{q\in\mathcal J_*}m_q(W;\varsigma).
}
\]
Then
\[
\operatorname{Cert}_{\mathfrak S}(W)=1
\quad\Longleftrightarrow\quad
r_{\mathfrak S}(W)>0.
\]

\begin{definition}[Specification-specific, robust, and fragile claims]
A claim is:
\begin{description}
\item[Specification-specific] when it is asserted only for one declared
 \(\varsigma_0\).
\item[Specification-robust] when every required gate has positive lower
 envelope across \(\mathfrak S\).
\item[Specification-fragile] when it passes for at least one admissible
 specification and fails for at least one other.
\end{description}
\end{definition}

A high fraction of passing specifications is informative but does not replace
the lower envelope when the claim is universal over \(\mathfrak S\). A family
in which ninety-nine analyses pass and one scientifically decisive hidden-
controller null fails does not support the complete relational certificate.
The noncompensatory logic remains in force at the specification level.

\subtopic{62.1 Geometry of a specification family}

Let \(d_{\mathfrak S}\) be a declared extended pseudometric on the quotient of
\(\mathfrak S\) by specification equivalence. A convenient decomposition is
\[
\begin{aligned}
d_{\mathfrak S}(\varsigma,\varsigma')
={}&
 w_Dd_D(\mathcal D,\mathcal D')
+w_Yd_Y(Y^+,Y^{+'})
+w_\ell d_\ell(\ell,\ell')\\
&+w_\tau d_\tau(\tau,\tau')
+w_Id_I(\mathcal I,\mathcal I')
+w_Nd_N(\mathcal N,\mathcal N')
+w_\Theta d_\Theta(\Theta,\Theta').
\end{aligned}
\]
No universal metric is imposed. The target term may use a divergence between
conditional laws under a declared coupling. The intervention and null terms
may use Hausdorff distances between implemented kernels. Threshold terms must
be expressed on normalized signed-margin scales. The metric is part of the
audit and should be varied in sensitivity analysis when more than one
scientifically defensible choice exists.

For a reference specification \(\varsigma_0\), define the distance to
certificate failure by
\[
\boxed{
r_{\mathrm{spec}}(W;\varsigma_0)
=
\inf\left\{
 d_{\mathfrak S}(\varsigma,\varsigma_0):
 \min_q m_q(W;\varsigma)\leq0
\right\}.
}
\]
If no failing specification exists in the declared universe, the radius is
infinite relative to that universe. A radius near zero means that a small
scientifically defensible change can reverse the conclusion.

\subtopic{62.2 Specification transport}

Changing the decomposition can alter the coordinate realization of the
workspace, while changing the target or horizon can alter the relevant
predictive state. Let
\[
\Phi_{\varsigma}:W_{\varsigma_0}\longrightarrow W_{\varsigma}
\]
be the declared correspondence when such transport is needed. For each gate,
assume a bound
\[
\left|
 m_q(\Phi_{\varsigma}W;\varsigma)
 -m_q(W;\varsigma_0)
\right|
\leq
L_qd_{\mathfrak S}(\varsigma,\varsigma_0)
+b_q(\varsigma),
\]
where \(b_q\) records mismatch not represented by the metric. The mismatch may
include a change of model class, imperfect target coupling, or an intervention
whose implementation cannot be transported exactly.

\subtopic{62.3 \(T_{28}\): Specification-transport robustness theorem}

\begin{theorem}[Specification-transport robustness]
\label{thm:specification-transport}
Let \(\mathfrak S\) be an admissible specification family with reference
\(\varsigma_0\). Suppose every required gate satisfies the transport bound
above. If
\[
\boxed{
 m_q(W;\varsigma_0)
>
\sup_{\varsigma\in\mathfrak S}
\left[
L_qd_{\mathfrak S}(\varsigma,\varsigma_0)
+b_q(\varsigma)
\right]
}
\]
for every \(q\in\mathcal J_*\), then
\[
\operatorname{Cert}(\Phi_{\varsigma}W;\varsigma)=1
\qquad
\text{for every }\varsigma\in\mathfrak S.
\]
If \(b_q\equiv0\), then
\[
\boxed{
r_{\mathrm{spec}}(W;\varsigma_0)
\geq
\min_{q:L_q>0}
\frac{m_q(W;\varsigma_0)}{L_q}.
}
\]
\end{theorem}

\begin{proof}
For each \(\varsigma\in\mathfrak S\),
\[
m_q(\Phi_{\varsigma}W;\varsigma)
\geq
m_q(W;\varsigma_0)
-L_qd_{\mathfrak S}(\varsigma,\varsigma_0)
-b_q(\varsigma).
\]
The hypothesis makes the right-hand side positive uniformly over the family.
Every required gate therefore passes under every admissible specification.
When \(b_q=0\), no gate can cross zero at a distance smaller than its reference
margin divided by its Lipschitz constant. Taking the minimum over gates gives
the radius bound.
\end{proof}

The theorem is deliberately gatewise. A large structural margin cannot absorb
a fragile causal margin. It also clarifies the relation to Part~XVII. The
decomposition-transport theorem is recovered by varying only the
\(\mathcal D\) coordinate while holding the remaining coordinates fixed.

\subtopic{62.4 A three-budget sufficient condition}

In practice, the reference margin is reduced by three distinct uncertainty
budgets:
\[
\boxed{
B_q^{\mathrm{total}}
=
B_q^{\mathrm{model}}
+B_q^{\mathrm{spec}}
+B_q^{\mathrm{sample}}.
}
\]
The first covers variation across observationally compatible realizations. The
second covers variation across admissible specifications. The third covers
finite-sample estimation. The sufficient condition
\[
\widehat m_q(W;\varsigma_0)
>
B_q^{\mathrm{total}}
\qquad(q\in\mathcal J_*)
\]
will reappear below in identified-set and confidence-envelope form. Keeping the
budgets separate prevents an apparently narrow confidence interval from being
mistaken for strong identification or specification robustness.

\topic{63}{Experimental Equivalence and Identifiability}

A complete mathematical definition does not guarantee that its truth can be
recovered from available data. The experimental design may fail to distinguish
two realizations that have different internal organization.

\begin{definition}[Experimental design and induced law]
\label{def:workspace-experiment-design}
An experimental design is a tuple
\[
\mathcal E
=
(\mathcal O,\mathcal I,\mathcal T),
\]
where \(\mathcal O\) is the observation map, \(\mathcal I\) is the family of
implemented interventions, and \(\mathcal T\) is the set of measured temporal
windows. A realization \(\mathfrak R\) induces the family of laws
\[
P_{\mathfrak R}^{\mathcal E}
=
\left
\{
\mathcal L_{\mathfrak R}
\bigl(\mathcal O(\mathbf H)\mid\operatorname{do}(\iota)\bigr):
\iota\in\mathcal I,
\ t\in\mathcal T
\right\}.
\]
\end{definition}

\begin{definition}[Experimental equivalence]
Two realizations are experimentally equivalent under \(\mathcal E\), written
\[
\mathfrak R\equiv_{\mathcal E}\mathfrak R',
\]
when
\[
P_{\mathfrak R}^{\mathcal E}
=
P_{\mathfrak R'}^{\mathcal E}.
\]
A property \(\psi(\mathfrak R)\) is identifiable on a model class
\(\mathfrak M\) under \(\mathcal E\) when it is constant on every
\(\equiv_{\mathcal E}\)-equivalence class contained in \(\mathfrak M\).
\end{definition}

Identifiability is therefore relative to both a model class and an experiment
family. Enlarging the intervention family refines the equivalence classes and
can make previously hidden distinctions identifiable. Enlarging the model
class can have the opposite effect.

\subtopic{63.1 \(T_{29}\): Non-identifiability impossibility theorem}

\begin{theorem}[No uniformly correct decision under experimental equivalence]
\label{thm:workspace-nonidentifiability}
Let \(\mathfrak R,\mathfrak R'\in\mathfrak M\) satisfy
\[
\mathfrak R\equiv_{\mathcal E}\mathfrak R'
\]
but
\[
\operatorname{Cert}(W_{\mathfrak R};\varsigma)
\neq
\operatorname{Cert}(W_{\mathfrak R'};\varsigma).
\]
Then the workspace certificate is not identifiable on \(\mathfrak M\) under
\(\mathcal E\). More strongly, for every possibly randomized decision rule
\(\widehat C\) based only on data generated by \(\mathcal E\),
\[
\boxed{
\mathbb P_{\mathfrak R}
\left(
\widehat C\neq C_{\mathfrak R}
\right)
+
\mathbb P_{\mathfrak R'}
\left(
\widehat C\neq C_{\mathfrak R'}
\right)
=1,
}
\]
where \(C_{\mathfrak R}=\operatorname{Cert}(W_{\mathfrak R};\varsigma)\).
Consequently,
\[
\inf_{\widehat C}
\max\left\{
\mathbb P_{\mathfrak R}(\widehat C\neq C_{\mathfrak R}),
\mathbb P_{\mathfrak R'}(\widehat C\neq C_{\mathfrak R'})
\right\}
\geq\frac12.
\]
\end{theorem}

\begin{proof}
The data law is identical under the two realizations, so
\(p=\mathbb P(\widehat C=1)\) is the same in both. If
\(C_{\mathfrak R}=0\) and \(C_{\mathfrak R'}=1\), the two error probabilities
are \(p\) and \(1-p\). Their sum is one, and at least one is at least one half.
The opposite labeling is identical.
\end{proof}

The theorem is not a limitation of a particular estimator. It is an
information-theoretic obstruction. No amount of optimization or sample size
can resolve a difference that the declared observations and interventions do
not expose.

\subtopic{63.2 Hidden-controller witness}

Consider two realizations that emit the same local streams \(X_1,X_2\) and the
same outcome
\[
Y=X_1\oplus X_2.
\]
In the distributed realization, the two streams are generated and retained by
separately intervenable carriers, and no carrier contains the pair. In the
centralized realization, a hidden controller contains \((X_1,X_2)\) and emits
the two displayed streams. Under passive observation of
\((X_1,X_2,Y)\), the induced laws are identical. Even interventions that alter
only the displayed messages may fail to distinguish them if the hidden
controller simply re-emits the manipulated values.

Carrier-relative relationality becomes identifiable only when the experiment
can test the claimed causal separation. Examples include independently
interrupting the candidate carriers, exposing the hidden memory, randomizing a
shared controller, or measuring whether one physical mechanism reconstructs
the complete state before the displayed split. The decomposition controls
\(N_{25}\)--\(N_{27}\) are therefore identifiability interventions, not merely
optional robustness checks.

\subtopic{63.3 Identifiability of individual gates}

The full certificate may be unidentified even when some coordinates are
identified. Structural dimension might be recoverable from observed transfer
constraints while causal specificity remains ambiguous. Conversely, a causal
contrast may be identifiable even when the latent workspace coordinates are
known only up to gauge. Each margin should therefore be assigned one of four
statuses:
\[
\text{point identified},
\quad
\text{sign identified},
\quad
\text{partially identified},
\quad
\text{unidentified}.
\]
A complete pass requires the sign of every required gate to be identified as
positive, not merely that one compatible model produces a positive estimate.

\topic{64}{Partial Identification of Workspace Claims}

When point identification fails, the correct response is not to choose one
convenient realization. It is to report the set of values compatible with the
data and assumptions \cite{Manski2003,ImbensManski2004}.

\begin{definition}[Identified realization set]
\label{def:identified-realization-set}
Given an experimental law \(P^{\mathcal E}\) and model class \(\mathfrak M\),
define
\[
\boxed{
\mathfrak R_{\mathrm{ID}}
(P^{\mathcal E};\mathfrak M)
=
\left\{
\mathfrak R\in\mathfrak M:
P_{\mathfrak R}^{\mathcal E}=P^{\mathcal E}
\right\}.
}
\]
For a gate \(q\) and specification \(\varsigma\), its identified margin set is
\[
\mathcal M_q^{\mathrm{ID}}
(P^{\mathcal E};\varsigma)
=
\left\{
 m_q(W_{\mathfrak R};\varsigma):
\mathfrak R\in\mathfrak R_{\mathrm{ID}}
\right\}.
\]
\end{definition}

The sharp lower and upper identified margins are
\[
\underline m_q^{\mathrm{ID}}(\varsigma)
=
\inf\mathcal M_q^{\mathrm{ID}}(\varsigma),
\qquad
\overline m_q^{\mathrm{ID}}(\varsigma)
=
\sup\mathcal M_q^{\mathrm{ID}}(\varsigma).
\]
A gate is sign-identified as passing if its lower bound is positive. It is
sign-identified as failing if its upper bound is nonpositive. If the interval
crosses zero, the available design does not determine that gate's outcome.

\subtopic{64.1 Identified certificate set}

Define
\[
\boxed{
\mathcal C_{\mathrm{ID}}
(P^{\mathcal E};\varsigma)
=
\left\{
\operatorname{Cert}(W_{\mathfrak R};\varsigma):
\mathfrak R\in\mathfrak R_{\mathrm{ID}}
\right\}
\subseteq\{0,1\}.
}
\]
The certificate is point identified exactly when this set is a singleton. The
marginal gate intervals provide simple sufficient conditions, but they need
not be logically necessary because gate margins can be dependent across the
identified set.

For robustness across specifications, define the realization-level robust
margin
\[
r_{\mathfrak S}(\mathfrak R)
=
\inf_{\varsigma\in\mathfrak S}
\min_{q\in\mathcal J_*}
 m_q(W_{\mathfrak R};\varsigma)
\]
and its identified range
\[
\underline r_{\mathrm{ID},\mathfrak S}
=
\inf_{\mathfrak R\in\mathfrak R_{\mathrm{ID}}}
 r_{\mathfrak S}(\mathfrak R),
\qquad
\overline r_{\mathrm{ID},\mathfrak S}
=
\sup_{\mathfrak R\in\mathfrak R_{\mathrm{ID}}}
 r_{\mathfrak S}(\mathfrak R).
\]

\subtopic{64.2 \(T_{30}\): Joint specification-identification theorem}

\begin{theorem}[Identification of a specification-robust certificate]
\label{thm:joint-spec-identification}
Let \(\mathfrak R_{\mathrm{ID}}\) be the identified realization set and
\(\mathfrak S\) an admissible specification family.
\begin{enumerate}
\item If
\[
\boxed{
\underline r_{\mathrm{ID},\mathfrak S}>0,
}
\]
then every realization compatible with the experimental law passes every
required gate under every admissible specification. The robust certificate is
identified as passing.
\item If
\[
\boxed{
\overline r_{\mathrm{ID},\mathfrak S}\leq0,
}
\]
then no compatible realization passes the complete certificate throughout the
family. The robust certificate is identified as failing.
\item If there exist compatible realizations
\(\mathfrak R_+,\mathfrak R_-\) with
\[
r_{\mathfrak S}(\mathfrak R_+)>0,
\qquad
r_{\mathfrak S}(\mathfrak R_-)\leq0,
\]
then the robust certificate is not identified under the declared experiment.
\end{enumerate}
\end{theorem}

\begin{proof}
The first two conclusions follow directly from the definitions of the infimum
and supremum over the identified set. The third gives two experimentally
compatible realizations with different robust-certificate values, so
Theorem~\ref{thm:workspace-nonidentifiability} applies.
\end{proof}

This theorem joins two forms of robustness that are often conflated. A result
can be stable across specifications within one fitted model yet fail to be
identified across models compatible with the data. It can also be identified
under one specification while changing sign across a reasonable specification
family. The strongest claim controls both axes.

\subtopic{64.3 Confidence regions for identified sets}

In finite samples, the identified set is itself uncertain. Let
\[
\widehat{\mathfrak R}_{1-\delta}
\]
be a confidence region satisfying
\[
\inf_{\mathfrak R_0\in\mathfrak M}
\mathbb P_{\mathfrak R_0}
\left(
\mathfrak R_0\in
\widehat{\mathfrak R}_{1-\delta}
\right)
\geq1-\delta.
\]
A conservative robust lower margin is
\[
\widehat r_{1-\delta}
=
\inf_{\mathfrak R\in\widehat{\mathfrak R}_{1-\delta}}
\inf_{\varsigma\in\mathfrak S}
\min_q m_q(W_{\mathfrak R};\varsigma).
\]
If \(\widehat r_{1-\delta}>0\), a false specification-robust certificate has
probability at most \(\delta\), provided that the model class and
specification family were fixed before the confidence region was computed.
This confidence-set approach is useful when gate estimators are strongly
coupled or when the latent realization is the primary inferential object.

\topic{65}{Uniform Finite-Sample Certification}

Population margins are not observed directly. A certification rule must
control the estimation error uniformly across every specification and gate on
which its conclusion depends.

Let \(\widehat m_q(\varsigma)\) estimate
\(m_q(W;\varsigma)\). Suppose there are nonnegative error radii
\(\varepsilon_q(n,\delta)\) such that
\[
\mathbb P\left(
\sup_{\varsigma\in\mathfrak S}
\left|
\widehat m_q(\varsigma)-m_q(\varsigma)
\right|
\leq
\varepsilon_q(n,\delta)
\quad\text{for all }q\in\mathcal J_*
\right)
\geq1-\delta.
\]
Define the empirical lower confidence envelope
\[
\boxed{
\widehat{\underline m}_{q,1-\delta}^{\mathfrak S}
=
\inf_{\varsigma\in\mathfrak S}
\widehat m_q(\varsigma)
-
\varepsilon_q(n,\delta).
}
\]

\subtopic{65.1 \(T_{31}\): Uniform specification-certificate theorem}

\begin{theorem}[Uniform finite-sample certification across specifications]
\label{thm:uniform-spec-certification}
Under the uniform error event above, declare a specification-robust certificate
only when
\[
\widehat{\underline m}_{q,1-\delta}^{\mathfrak S}>0
\qquad
\text{for every }q\in\mathcal J_*.
\]
Then
\[
\boxed{
\sup_{P:\,\operatorname{Cert}_{\mathfrak S}(W)=0}
\mathbb P_P
\left(
\text{the robust certificate is declared}
\right)
\leq\delta.
}
\]
\end{theorem}

\begin{proof}
On the uniform error event,
\[
m_q(\varsigma)
\geq
\widehat m_q(\varsigma)-\varepsilon_q
\]
for every gate and every specification. If all empirical lower envelopes are
positive, every population margin is positive throughout the family. A false
declaration can therefore occur only when the uniform error event fails, which
has probability at most \(\delta\).
\end{proof}

The theorem controls the complete specification family, not merely the
analysis that happened to produce the strongest estimate. It is an extension
of the gatewise intersection rule in Theorem~\ref{thm:gatewise-validity}.
Gatewise validity controls the conjunction at one frozen specification.
Uniform validity controls the conjunction over a family.

\subtopic{65.2 Finite specification grids}

Suppose \(\mathfrak S\) and \(\mathcal J_*\) are finite, and each margin is
estimated as
\[
\widehat m_q(\varsigma)
=
\frac1n\sum_{i=1}^nZ_{i,q,\varsigma}-\theta_{q,\varsigma},
\qquad
0\leq Z_{i,q,\varsigma}\leq1,
\]
from independent certification units. Hoeffding's inequality and a union bound
give
\[
\mathbb P\left(
\sup_{q,\varsigma}
\left|
\widehat m_q(\varsigma)-m_q(\varsigma)
\right|
>
\sqrt{
\frac{
\log\bigl(2|\mathcal J_*||\mathfrak S|/\delta\bigr)
}{2n}}
\right)
\leq\delta.
\]
Thus the choice
\[
\boxed{
\varepsilon_n(\delta)
=
\sqrt{
\frac{
\log\bigl(2|\mathcal J_*||\mathfrak S|/\delta\bigr)
}{2n}
}
}
\]
yields an explicit conservative certificate. Dependent episodes require a
bound appropriate to the true independent unit rather than token-level or
edge-level pseudo-replication.

\subtopic{65.3 Continuous specification families}

Let \(\mathfrak S_\eta\) be an \(\eta\)-net of
\((\mathfrak S,d_{\mathfrak S})\) with covering number
\(N(\eta,\mathfrak S,d_{\mathfrak S})\). If both the population and empirical
margin functions are \(L_q\)-Lipschitz, then
\[
\sup_{\varsigma\in\mathfrak S}
|\widehat m_q(\varsigma)-m_q(\varsigma)|
\leq
\max_{\varsigma_j\in\mathfrak S_\eta}
|\widehat m_q(\varsigma_j)-m_q(\varsigma_j)|
+2L_q\eta.
\]
For bounded empirical means, a valid radius is therefore
\[
\boxed{
\varepsilon_{q,n,\eta}(\delta)
=
\sqrt{
\frac{
\log\bigl(2|\mathcal J_*|
N(\eta,\mathfrak S,d_{\mathfrak S})/\delta\bigr)
}{2n}
}
+2L_q\eta.
}
\]
The tuning parameter \(\eta\) balances net cardinality against interpolation
error. More refined empirical-process bounds may replace this elementary
covering argument without changing the certificate logic.

\subtopic{65.4 Power and indifference regions}

A robust procedure can be valid yet uninformative when the family is too large
or the data are insufficient. Each gate should therefore report an
indifference region \([-\kappa_q,\kappa_q]\). A conclusion is:
\begin{description}
\item[Robust pass] when the lower confidence envelope exceeds
 \(\kappa_q\) for every gate.
\item[Robust fail] when a required gate's upper confidence envelope is below
 \(-\kappa_q\) throughout the relevant family.
\item[Specification conflict] when well-estimated admissible specifications
 lie on opposite sides of the indifference region.
\item[Statistically unresolved] when sampling uncertainty intersects the
 indifference region.
\end{description}
This vocabulary separates genuine sensitivity from low power.

\topic{66}{Adaptive Specification Choice and Selective Validity}

A pointwise valid test does not remain valid after the analyst searches for the
most favorable target, horizon, null panel, or threshold. Consider
\(K\) independent null \(p\)-values and select the smallest. Then
\[
\mathbb P\left(\min_{1\leq k\leq K}p_k\leq\alpha\right)
=
1-(1-\alpha)^K,
\]
which exceeds \(\alpha\) whenever \(K>1\). The same logic applies when a
workspace certificate is reported only for the specification that produced the
largest gate margins.

The issue is not exploration itself. Exploration is often necessary to learn
which target and intervention are scientifically meaningful. The issue is
using the same random outcomes both to choose the claim and to certify it.
Post-selection inference addresses the general problem by using simultaneous
or selection-adjusted procedures \cite{BerkEtAl2013}.

\subtopic{66.1 Specification leakage}

\begin{definition}[Specification leakage]
Specification leakage occurs when certification outcomes influence any of the
following without a fresh inferential correction:
\[
Y^+,
\ell,
\tau,
\mathcal I,
\mathcal N,
\Theta,
\mathfrak S,
\mathcal J_*.
\]
The leakage status is separate from formation leakage. A workspace extractor
may be frozen while the analyst still searches over targets or thresholds on
the certification split.
\end{definition}

A valid report must retain the complete search history, including failed and
abandoned specifications. Reporting only the selected analysis prevents the
reader from determining the effective multiplicity.

\subtopic{66.2 \(T_{32}\): Sample-split selection theorem}

\begin{theorem}[Validity after independent specification selection]
\label{thm:sample-split-spec-selection}
Let \(\widehat\varsigma\) be selected using
\(\mathcal D_{\mathrm{sel}}\), and let
\(\mathcal D_{\mathrm{cert}}\) be independent of
\(\mathcal D_{\mathrm{sel}}\). For every fixed
\(\varsigma\in\mathfrak S\), let
\(\phi_\varsigma(\mathcal D_{\mathrm{cert}})\in\{0,1\}\) be a certificate
rule satisfying
\[
\sup_{P:\operatorname{Cert}(W;\varsigma)=0}
\mathbb P_P\bigl(\phi_\varsigma=1\bigr)
\leq\alpha.
\]
Then the selected-specification procedure satisfies
\[
\boxed{
\sup_P
\mathbb P_P\left(
\phi_{\widehat\varsigma}=1,
\operatorname{Cert}(W;\widehat\varsigma)=0
\right)
\leq\alpha.
}
\]
\end{theorem}

\begin{proof}
Condition on \(\mathcal D_{\mathrm{sel}}\). The selected specification is then
fixed. Independence leaves the certification data distributed exactly as in
the fixed-specification test, whose false-pass probability is at most
\(\alpha\). Averaging the conditional bound over the selection data preserves
the same bound.
\end{proof}

The theorem justifies discovery on formation or calibration data followed by a
single frozen evaluation on fresh certification data. Reusing the
certification split to revise the specification requires either another fresh
split or a simultaneous post-selection procedure whose validity covers the
entire search.

\subtopic{66.3 Nested selection and stress testing}

A practical design may use:
\[
\mathcal D_{\mathrm{form}}
\sqcup
\mathcal D_{\mathrm{spec}}
\sqcup
\mathcal D_{\mathrm{cert}}
\sqcup
\mathcal D_{\mathrm{stress}}.
\]
Formation learns the workspace. The specification split chooses targets,
horizons, nulls, and thresholds from a predeclared menu. Certification makes
the pass decision once. Stress testing explores broader shifts without
retroactively changing the certified claim. Cross-fitting can recover sample
efficiency, but the fold-level selection and evaluation paths must remain
separated and the final uncertainty calculation must account for recombination.

\subtopic{66.4 Thresholds and null families are part of selection}

Threshold tuning can be as consequential as target selection. If the observed
causal effect is \(0.12\), choosing a minimum effect threshold of \(0.10\)
after seeing the estimate is not certification. Likewise, retaining only the
nulls that the intact system defeats converts a falsification panel into a
success filter. Thresholds and null families must be frozen or evaluated under
a uniform rule that includes every considered alternative.

\topic{67}{Specification Audit and Reporting}

A complete specification audit should accompany every strong relational
workspace claim.

\subtopic{67.1 Audit procedure}

\begin{enumerate}
\item State the reference specification \(\varsigma_0\) in full.
\item Distinguish equivalent recodings, robustness alternatives, and genuinely
 different scope claims.
\item Declare the admissible family \(\mathfrak S\) and justify every major
 exclusion before certification.
\item Record the specification pseudometric or other transport bounds used to
 compare family members.
\item Compute every gate's lower envelope, upper envelope, sensitivity width,
 and reference margin.
\item State the experiment family \(\mathcal E\) and the model class
 \(\mathfrak M\) relative to which identifiability is claimed.
\item Report identified sets or bounds whenever point identification fails.
\item Apply a uniform finite-sample rule over every gate and specification used
 in the conclusion.
\item Preserve the complete specification-selection history and the data split
 on which each choice was made.
\item Report the strongest pointwise, family-robust, sign-identified, and
 statistically resolved conclusions separately.
\end{enumerate}

\subtopic{67.2 Additional adversarial controls}

The empirical null and audit panel is extended by:
\[
\begin{aligned}
N_{29}&=\text{post hoc target-selection audit},\\
N_{30}&=\text{horizon cherry-picking audit},\\
N_{31}&=\text{loss and utility substitution audit},\\
N_{32}&=\text{restricted intervention-family audit},\\
N_{33}&=\text{threshold and gate-set tuning audit},\\
N_{34}&=\text{frozen specification sweep across }\mathfrak S.
\end{aligned}
\]
The first five reconstruct or bound the analysis search that preceded the
reported claim. The final control reruns the complete frozen certificate across
the admissible family. A claim that survives \(N_{34}\) but fails the search
audit remains vulnerable to selection bias. A claim that passes the search
audit but changes sign across \(\mathfrak S\) is validly estimated yet
specification-fragile.

\subtopic{67.3 Required specification record}

The machine-readable record should include:
\[
\boxed{
\begin{aligned}
\mathsf{SpecificationRecord}
=\{&
\mathsf{reference},
\mathsf{family},
\mathsf{equivalence\_classes},
\mathsf{metric},\\
&\mathsf{targets},
\mathsf{losses},
\mathsf{horizons},
\mathsf{interventions},\\
&\mathsf{null\_panels},
\mathsf{thresholds},
\mathsf{selection\_history},\\
&\mathsf{identified\_sets},
\mathsf{uniform\_bounds},
\mathsf{gate\_envelopes}
\}.
\end{aligned}
}
\]
Every omitted family member requires a reason. Every changed specification
requires a data-split label and time stamp. Every gate requires pointwise and
family-wide status.

\subtopic{67.4 Claim-strength vocabulary}

The following vocabulary is recommended:
\begin{description}
\item[Pointwise demonstrated.] The certificate passes at one frozen
 specification with valid uncertainty.
\item[Specification-robust.] It passes uniformly across a declared admissible
 family.
\item[Sign identified.] Every realization compatible with the experimental
 law gives the same gate sign.
\item[Robustly identified.] The complete family-robust certificate is
 identified as passing.
\item[Partially identified.] The data restrict the certificate margins but do
 not determine every required sign.
\item[Specification-fragile.] Validly estimated admissible specifications give
 different certificate outcomes.
\item[Selection-contaminated.] The reported specification was chosen using
 certification outcomes without valid adjustment.
\end{description}

\subtopic{67.5 Consequence for the motivating collaboration}

The motivating human--AI research process remains a plausible candidate
persistent workspace, but its strongest current description is narrower than
a certified robust claim. The relevant target trajectory, horizon family,
intervention set, null panel, thresholds, and admissible specification family
have not been frozen and tested prospectively. Nor has the carrier-relative
workspace predicate been shown to be identifiable against hidden-controller
and matched-storage alternatives. The theory therefore supports a disciplined
research program rather than a retrospective declaration.

The terminal architecture is now explicit. A strong claim must survive three
separate questions:
\[
\boxed{
\begin{aligned}
&\text{Would it survive another defensible specification?}\\
&\text{Could a contrary realization generate the same experiments?}\\
&\text{Does the finite sample support every required margin uniformly?}
\end{aligned}
}
\]
Once these questions are answered, further extensions concern new mathematical
settings or new empirical domains rather than a missing foundational layer.
